% GENERATED FILE. Edit canonical lessons, then run npm run generate:books.
% canonical-source-hash: fnv1a64:48d45e441ef7278d
% canonical-lessons: LA-C452-aberro, LA-C452-devius, LA-C452-invius, LA-C452-ambages, LA-C452-vagor

\chapter{To stray, Off the road, Without roads, Roundabout way, To wander about}
\label{ch:la-to-stray-off-the-road-without-roads-roundabout-way-to-wander-about}
\hlchaptermodality{452}

\begin{chapteropening}
\textbf{By the end of this chapter:} \emph{I can talk about getting lost: straying from the road, going off the road, roadless country, a roundabout way, and wandering about.}

Complete the last lesson of chapter 452: I can talk about getting lost: straying from the road, going off the road, roadless country, a roundabout way, and wandering about.
\end{chapteropening}

\section[aberrō, aberrāre]{aberrō, aberrāre — to stray, to wander from the way}

\label{lesson:LA-C452-aberro}



\begin{quote}
Before the new one: say the Latin for to show, then the Latin for to make known.
\end{quote}

\subsection*{You'll want to know: aberrō, aberrāre}
\textbf{aberrō, aberrāre} — \textquotedblleft{}to stray, to wander from the way\textquotedblright{}.

\begin{quote}
In nebulā viātōrēs ā viā aberrant.
\end{quote}

\begin{quote}
\emph{In the fog the travellers stray from the road.}
\end{quote}

\emph{aberrō, aberrāre, aberrāvī, aberrātum}, with \emph{ā} or \emph{ab}. English \emph{aberration} comes from it.

\begin{cousinweb}
\emph{aberrō} is built from \emph{ab} 'away, from' and \emph{errō} 'to wander, to go astray' (EtymDB; Lewis \& Short). English \emph{aberration} comes from it.
\end{cousinweb}

\begin{grammarlens}[title={What you've built}]
One more word for telling, explaining and describing at B1.
\end{grammarlens}

\subsection*{Your turn}
\begin{itemize}
\raggedright
  \item \emph{Say it:} \emph{aberrō, aberrāre}
  \item \emph{Say it:} \emph{aberrō, aberrāre}, once more
  \item \emph{Say it:} a sentence of your own with \emph{aberrō, aberrāre}
  \item \emph{Recall:} say the Latin for to show, then the Latin for to make known, then say \emph{aberrō, aberrāre} again
\end{itemize}

\begin{tcolorbox}[breakable,colback=teal!4,colframe=teal!35!black,title={Before you move on}]
What does aberrō, aberrāre mean? (\textquotedblleft{}To stray, to wander from the way\textquotedblright{}.) Say it once more.
\end{tcolorbox}
\section[dēvius, dēvia, dēvium]{dēvius, dēvia, dēvium — off the road, out of the way; remote}

\label{lesson:LA-C452-devius}



\begin{quote}
Before the new one: say the Latin for to make known, then the Latin for to stray.
\end{quote}

\subsection*{You'll want to know: dēvius, dēvia, dēvium}
\textbf{dēvius, dēvia, dēvium} — \textquotedblleft{}off the road, out of the way; remote\textquotedblright{}.

\begin{quote}
Vīlla dēvia est; nēmō eam facile invenit.
\end{quote}

\begin{quote}
\emph{The country house is out of the way; nobody finds it easily.}
\end{quote}

An adjective of three endings, from \emph{dē} 'off' and \emph{via} 'road'. English \emph{devious} first meant 'off the main road'.

\begin{cousinweb}
\emph{dēvius} is built from \emph{dē} 'off, away from' and \emph{via} 'road' (Lewis \& Short). English \emph{devious} comes from it, first in the sense 'out of the way'.
\end{cousinweb}

\begin{grammarlens}[title={What you've built}]
One more word for telling, explaining and describing at B1.
\end{grammarlens}

\subsection*{Your turn}
\begin{itemize}
\raggedright
  \item \emph{Say it:} \emph{dēvius, dēvia, dēvium}
  \item \emph{Say it:} \emph{dēvius, dēvia, dēvium}, once more
  \item \emph{Say it:} a sentence of your own with \emph{dēvius, dēvia, dēvium}
  \item \emph{Recall:} say the Latin for to make known, then the Latin for to stray, then say \emph{dēvius, dēvia, dēvium} again
\end{itemize}

\begin{tcolorbox}[breakable,colback=teal!4,colframe=teal!35!black,title={Before you move on}]
What does dēvius, dēvia, dēvium mean? (\textquotedblleft{}Off the road, out of the way; remote\textquotedblright{}.) Say it once more.
\end{tcolorbox}
\section[invius, invia, invium]{invius, invia, invium — without roads, pathless, impassable}

\label{lesson:LA-C452-invius}



\begin{quote}
Before the new one: say the Latin for to stray, then the Latin for off the road.
\end{quote}

\subsection*{You'll want to know: invius, invia, invium}
\textbf{invius, invia, invium} — \textquotedblleft{}without roads, pathless, impassable\textquotedblright{}.

\begin{quote}
Silva magna et invia est; equī trānsīre nōn possunt.
\end{quote}

\begin{quote}
\emph{The wood is big and pathless; the horses cannot get through.}
\end{quote}

An adjective of three endings, from \emph{in-} 'not' and \emph{via} 'road'.

\begin{cousinweb}
\emph{invius} is built from \emph{in-} 'not' and \emph{via} 'road' (EtymDB; Lewis \& Short): without a road, so hard to pass.
\end{cousinweb}

\begin{grammarlens}[title={What you've built}]
One more word for telling, explaining and describing at B1.
\end{grammarlens}

\subsection*{Your turn}
\begin{itemize}
\raggedright
  \item \emph{Say it:} \emph{invius, invia, invium}
  \item \emph{Say it:} \emph{invius, invia, invium}, once more
  \item \emph{Say it:} a sentence of your own with \emph{invius, invia, invium}
  \item \emph{Recall:} say the Latin for to stray, then the Latin for off the road, then say \emph{invius, invia, invium} again
\end{itemize}

\begin{tcolorbox}[breakable,colback=teal!4,colframe=teal!35!black,title={Before you move on}]
What does invius, invia, invium mean? (\textquotedblleft{}Without roads, pathless, impassable\textquotedblright{}.) Say it once more.
\end{tcolorbox}
\section[ambāgēs, ambāgum]{ambāgēs, ambāgum — a roundabout way, a winding path; long-winded talk}

\label{lesson:LA-C452-ambages}



\begin{quote}
Before the new one: say the Latin for off the road, then the Latin for without roads.
\end{quote}

\subsection*{You'll want to know: ambāgēs, ambāgum}
\textbf{ambāgēs, ambāgum} — \textquotedblleft{}a roundabout way, a winding path; long-winded talk\textquotedblright{}.

\begin{quote}
Pastor nōn rēctā viā sed per ambāgēs nōs dūcit.
\end{quote}

\begin{quote}
\emph{The shepherd leads us not by the direct road but by a roundabout way.}
\end{quote}

Third declension, feminine, mostly plural: \emph{ambāgēs, ambāgum f.}; the singular appears only in the ablative \emph{ambāge}. Also talk that goes round the point.

\begin{cousinweb}
\emph{ambāgēs} is built from \emph{ambi-} 'around' and \emph{agō} 'to drive, to go' (EtymDB): a going round. Lewis \& Short give a roundabout way, then roundabout speech.
\end{cousinweb}

\begin{grammarlens}[title={What you've built}]
One more word for telling, explaining and describing at B1.
\end{grammarlens}

\subsection*{Your turn}
\begin{itemize}
\raggedright
  \item \emph{Say it:} \emph{ambāgēs, ambāgum}
  \item \emph{Say it:} \emph{ambāgēs, ambāgum}, once more
  \item \emph{Say it:} a sentence of your own with \emph{ambāgēs, ambāgum}
  \item \emph{Recall:} say the Latin for off the road, then the Latin for without roads, then say \emph{ambāgēs, ambāgum} again
\end{itemize}

\begin{tcolorbox}[breakable,colback=teal!4,colframe=teal!35!black,title={Before you move on}]
What does ambāgēs, ambāgum mean? (\textquotedblleft{}A roundabout way, a winding path; long-winded talk\textquotedblright{}.) Say it once more.
\end{tcolorbox}
\section[vagor, vagārī]{vagor, vagārī — to wander about, to roam}

\label{lesson:LA-C452-vagor}



\begin{quote}
Before the new one: say the Latin for without roads, then the Latin for a roundabout way.
\end{quote}

\subsection*{You'll want to know: vagor, vagārī}
\textbf{vagor, vagārī} — \textquotedblleft{}to wander about, to roam\textquotedblright{}.

\begin{quote}
Canis sine dominō per viās oppidī vagātur.
\end{quote}

\begin{quote}
\emph{A dog without a master roams the streets of the town.}
\end{quote}

Deponent: \emph{vagor, vagārī, vagātus sum}. English \emph{vagabond} comes from \emph{vagābundus}, built on it.

\begin{cousinweb}
\emph{vagor} is built on the adjective \emph{vagus} 'wandering, roving' (EtymDB; Lewis \& Short). English \emph{vague} comes from \emph{vagus}, and \emph{vagabond} from \emph{vagābundus}, built on the verb.
\end{cousinweb}

\begin{grammarlens}[title={What you've built}]
One more word for telling, explaining and describing at B1.
\end{grammarlens}

\subsection*{Your turn}
\begin{itemize}
\raggedright
  \item \emph{Say it:} \emph{vagor, vagārī}
  \item \emph{Say it:} \emph{vagor, vagārī}, once more
  \item \emph{Say it:} a sentence of your own with \emph{vagor, vagārī}
  \item \emph{Recall:} say the Latin for without roads, then the Latin for a roundabout way, then say \emph{vagor, vagārī} again
\end{itemize}

\begin{tcolorbox}[breakable,colback=teal!4,colframe=teal!35!black,title={Before you move on}]
What does vagor, vagārī mean? (\textquotedblleft{}To wander about, to roam\textquotedblright{}.) Say it once more.
\end{tcolorbox}
